% year: תשפ"ב
% type: מבחן
% semester: א
% moed: א
% number: 6א
% points: 8
% categories: integral-applications
% source: exams/מבחנים/תשפב/מועד א תשפב.pdf

\begin{question}
חשבו את השטח הסופי החסום בין הפונקציה $f(x)=x^2$ והישר $y=-x+2$.
\end{question}

\begin{hint}
מצאו את נקודות החיתוך, וקבעו איזו פונקציה עליונה בין שתיהן.
\end{hint}

\begin{solution}
\textbf{נקודות חיתוך:} $x^2=-x+2\iff x^2+x-2=0\iff (x+2)(x-1)=0$, כלומר $x=-2$ ו-$x=1$.
בקטע $(-2,1)$ מתקיים $-x+2-x^2=-(x+2)(x-1)>0$, ולכן הישר מעל הפרבולה. השטח:
\[ S=\int_{-2}^{1}\left(-x+2-x^2\right)dx=\left[-\frac{x^2}{2}+2x-\frac{x^3}{3}\right]_{-2}^{1}
=\left(-\frac12+2-\frac13\right)-\left(-2-4+\frac83\right)=\frac76+\frac{10}{3}=\boxed{\frac92}. \]
\end{solution}
